%%
%% spring-lecture05.tex
%% 
%% Made by Alex Nelson <pqnelson@gmail.com>
%% Login   <alex@lisp>
%% 
%% Started on  2026-04-09T11:06:05-0700
%% Last update 2026-04-09T11:06:05-0700
%% 

\lecture[]

\subsection{Length, volume, integration}

\begin{node}
When people speak of ``geometric properties'' (as opposed to
``topological properties''), we need to involve some ``measure'' (in
the ``measure theoretic'' sense). This
distinguishes geometric from topological properties. ``Measure''
invariably involves metrics.
\end{node}

\begin{definition}[Length of smooth curve]
Let $(M,g)$ be a Riemannian manifold. Let $\gamma\colon[a,b]\to M$ be
a smooth curve.\footnote{How can $\gamma$ be considered smooth on
$[a,b]$? This $[a,b]$ is not a manifold. The problem? The boundary!
The easiest way is that $\gamma$ can be extended to
$\widetilde{\gamma}\colon(a',b')\to M$ where $[a,b]\subset(a',b')$ and $\widetilde{\gamma}|_{[a,b]}=\gamma$.}
The \define{Length} of $\gamma$ is given by
\begin{equation}
\length_{g}(\gamma)=\int^{b}_{a}\sqrt{g(\gamma(t))[\gamma'(t),\gamma'(t)]}\,\D t.
\end{equation}
It is not uncommon to see this written as
$\int^{b}_{a}|\gamma'(t)|\,\D t$.
\end{definition}

\begin{definition}
Let $(M,g)$ be a Riemannian manifold. For every pair of points $x,y\in M$,
let
\begin{equation}
\mathcal{C}(x,y)=\{\gamma\colon[a,b]\to M\mbox{ smooth}\mid\gamma(a)=x,\gamma(b)=y\}.
\end{equation}
We define the \define{Riemannian Distance} between $x$ and $y$ by
\begin{equation}
d_{g}(x,y):=\inf_{\gamma\in\mathcal{C}(x,y)}\length_{g}(\gamma).
\end{equation}
(If $M$ is disconnected and $x$ and $y$ live on different components,
then $d_{g}(x,y)=+\infty$.)
\end{definition}

\begin{proposition}
If $M$ is connected, then $(M,d_{g})$ is a metric space.
\end{proposition}

\begin{xca}
Prove this!
\end{xca}

\begin{node}
The next notions are measure theoretic, but it is the most formal way
to describe these notions.
\end{node}

\begin{definition}[Volume measure]
Let $(M,g)$ be a Riemannian manifold. Let $\mathcal{B}(M)$ be the
Borel $\sigma$-algebra of $M$ (since $M$ is a topological space, we
can take the smallest $\sigma$-algebra containing the topology [i.e.,
  the Borel $\sigma$-algebra]).
Let $\{(U_{k},\varphi_{k})\}_{k\in\NN}\subset\smoothstr{M}$ be a
countable smooth atlas, and let $\{\rho_{k}\}_{k\in\NN}$ be a
locally-finite partition of unity subordinated to the open cover
$\{U_{k}\}_{k\in\NN}$. Then for any $A\in\mathcal{B}(M)$, we define
the
\begin{equation}
  \mu_{g}(A):=\sum_{k\in\NN}\int_{\varphi_{k}(A\cap U_{k})}(\rho_{k}\circ\varphi^{-1}_{k})\sqrt{\det(g^{(k)}_{ij})}\,\D\mathcal{L}^{n}
\end{equation}
The completion of $\mu_{g}$ is a Radon measure on $M$ called the
\define{Volume Measure} and is denoted by $\vol_{g}$.
\end{definition}

\begin{remark}
\begin{enumerate}
\item If $A\propersubset U_{k}$, then all other $U_{j}$ contributions
  vanish and $\mu_{g}(A)=\int_{\varphi_{k}(A\cap U_{k})}(\rho_{k}\circ\varphi^{-1}_{k})\sqrt{\det(g^{(k)}_{ij})}\,\D\mathcal{L}^{n}$
\item We pullback integration from $\RR^{n}$ to the manifold
\item $\varphi_{k}(A\cap U_{k})$ is a Borel set in $\RR^{n}$, so we
  can integrate on it.
\item We should check $\mu\colon\mathcal{B}(M)\to[0,+\infty]$ is a
  Radon measure on $M$ (suffices to prove $\mu$ is an outer measure,
  which implies $\mu$ is a Borel measure; we need to prove more for
  establishing it is a Radon measure, though).
\item For measure theoretic reasons, we will want to work with the
  completion of $\mu$ to define a measure on
  $\mathcal{B}(M)\cup\{\mbox{null sets}\}$.
\end{enumerate}
\end{remark}

\begin{remark}
Some people want to relate this to integration of differential
forms. For \emph{orientable} manifolds, the top-form may be identified
as the volume measure. Actually \emph{proving} this claim is quite
tedious.

However, for non-orientable manifolds, our definition is the only game
in town.
\end{remark}

\begin{definition}[Volume of regions]
Let $(M,g)$ be a Riemannian manifold. Let $D\subset M$ be a ``region''
(i.e., a relatively compact\footnote{A set is ``relatively compact''
if its closure is compact.} open subset of $M$). Then the
\define{Volume} of $D$ is
\begin{equation}
\vol(D):=\vol_{g}(D)=\int_{D}\D\vol_{g}<+\infty.
\end{equation}
The ``relative compactness'' condition ensures the integral is finite.
\end{definition}

\subsection{Connections on smooth manifolds}

\begin{node}
We can differentiate a function $f$ with respect to a vector field $X$
by taking its directional derivative. For us to extend this procedure
from a function $f$ to a vector field, we need something more: a connection.
\end{node}

\begin{definition}
Let $M$ be a smooth manifold. A \define{Connection} on $TM$ (or,
abusing language, ``on $M$'') is an $\RR$-linear operator
\begin{equation}
\begin{split}
\nabla\colon&\Vect{M}\times\Vect{M}\to\Vect{M}\\
&(X,Y)\mapsto\nabla_{X}Y
\end{split}
\end{equation}
(we are differentiating $Y$ with respect to $X$) such that
\begin{enumerate}
\item for all $f\in C^{\infty}(M)$, we have $\nabla_{fX}Y=f\nabla_{X}Y$
\item \textbf{Leibniz rule}: for all $f\in C^{\infty}(M)$, we have $\nabla_{X}(fY)=(Xf)Y+f\nabla_{X}Y$
\end{enumerate}
This is also called the \define{Covariant Derivative}.
\end{definition}

\begin{remark}
\textit{A priori} there is no canonical way to construct a connection:
we can have many. But there is an obvious way to do it on a Riemannian manifold.
\end{remark}

\begin{definition}[Torsion-free]
Let $M$ be a smooth manifold.
Let $\nabla$ be a connection on $M$.
We say $\nabla$ is \define{Torsion-Free} if for all vector fields
$X,Y\in\Vect{M}$ we have
\begin{equation}
\nabla_{X}Y-\nabla_{Y}X=[X,Y].
\end{equation}
This is because the \define{Torsion} $\tau$ is defined as a
$(0,2)$-tensor satisfying (up to an overall sign):
\begin{equation}
\tau(X,Y) = \nabla_{X}Y-\nabla_{Y}X - [X,Y]
\end{equation}
\end{definition}

\begin{example}
In $\RR^{n}$, vector fields look like $X=(\id,v)\colon\RR^{n}\to\RR^{n}\times\RR^{n}$
and $Y=(\id,w)\colon\RR^{n}\to\RR^{n}\times\RR^{n}$. Then if we define
\begin{equation}
\nabla_{X}^{0}Y=(\id,\sum^{n}_{i=1}v_{i}\partial_{i}w),
\end{equation}
then this is called the \define{Euclidean Connection} on $\RR^{n}$. It
is torsion-free. (Later, we will see it is also the Levi-Civita connection.)
\end{example}